% part: intuitionistic-logic
% chapter: propositions-as-types
% section: reduction

\documentclass[../../../include/open-logic-section]{subfiles}

\begin{document}

\olfileid[tr]{int}{pty}{red}

\olsection{İndirgeme}

Doğal tümdengelim !!{derivation}{}inde, bir !!a{introduction} kuralını izleyen
bir !!a{elimination} kuralı ``indirgenebilen'' bir dolambaçtır. Örneğin
\Intro{\land}'ı izleyen \Elim{\land} ``birbirini götürür''; çünkü
\Elim{\land}'ın sonucu her zaman \Intro{\land}'ın ürettiği ve bağlaçlı
yapının bileşenlerinden biridir ve \Intro{\land}'ın öncülleri de iki
bağlaşıktır. Dolayısıyla bağlaşıklardan herhangi birinin bir
!!a{derivation}{}ine gereksinimimiz varsa ve bağlacını tanıtıp sonra
gidermek yerine o !!{derivation}{}i doğrudan kullanabiliriz:
\begin{gather*}
  \bottomAlignProof
  \AxiomC{}
  \RightLabel{$\delta_1$}
  \DeduceC{$!A_1$}
  \AxiomC{}
  \RightLabel{$\delta_2$}
  \DeduceC{$!A_2$}
  \RightLabel{$\Intro{\land}$}
  \BinaryInfC{$!A_1 \land !A_2$}
  \RightLabel{$\Elim{\land}_i$}
  \UnaryInfC{$!A_i$}
  \DisplayProof\bottomAlignProof
  \quad
  \redone
  \quad
  \AxiomC{}
  \RightLabel{$\delta_i$}
  \DeduceC{$!A_i$}
  \DisplayProof
\end{gather*}

Bu sadeleştirmenin düşündürdüğü biçimde, ilişkili ispat terimleri üzerinde
benzer bir indirgeme bağıntısı ele alabiliriz. $N_1$,~$\delta_1$'in;
$N_2$ de~$\delta_2$'nin ispat terimiyse soldaki ispatın ispat teriminin
(bir adımda) sağdaki ispatın ispat terimine indirgendiğini söyleyebiliriz:
\[
  \proj{i}{\pair{N_1}{N_2}} \redone N_i
\]
Tipli lambda kalkülüsünde bu, çarpım tipi için \emph{beta indirgeme}
kuralıdır.

Doğal tümdengelimde \Intro{\land}'ı izleyen \Elim{\land} gibi, yani
indirgemeye konu olan bir çıkarım çiftine \emph{kesme} denir. Tipli lambda
kalkülüsünde $\redone$ simgesinin solundaki terime \emph{indirgenen},
sağındaki terime \emph{indirgenmiş} denir. Yalnız
$(\lambd[x][N])Q$'nun indirgenen sayıldığı tiplenmemiş lambda kalkülüsünden
farklı olarak tipli lambda kalkülüsünün sözdizimi $\pair{N}{M}$,
$\proj{i}{N}$, $\inj[!A]{i}{N}$, $\dcase{N}{x_1}{M_1}{x_2}{M_2}$ ve
$\abort{!A}{N}$ içeren terimlerle, bunlara karşılık gelen indirgenenleri de
kapsayacak biçimde genişletilmiştir.

Veya bağlacı için de benzer bir indirgeme vardır:
\begin{gather*}
  \bottomAlignProof
  \AxiomC{}
  \RightLabel{$\delta$}
  \DeduceC{$!A_i$}
  \RightLabel{$\Intro{\lor}$}
  \UnaryInfC{$!A_1 \lor !A_2$}
  \AxiomC{$\Discharge{!A_1}{x_1}$}
  \RightLabel{$\delta_1$}
  \DeduceC{$!C$}
  \AxiomC{$\Discharge{!A_2}{x_2}$}
  \RightLabel{$\delta_2$}
  \DeduceC{$!C$}
  \DischargeRule{$\Elim{\lor}$}{x_1,x_2}
  \TrinaryInfC{$!C$}
  \DisplayProof\bottomAlignProof
  \quad
  \redone
  \quad
  \AxiomC{}
  \RightLabel{$\delta$}
  \DeduceC{$!A_i$}
  \RightLabel{$\delta_i$}
  \DeduceC{$!C$}
  \DisplayProof
\end{gather*}
Bu, ispat terimleri üzerinde şu indirgemeye karşılık gelir:
\begin{gather*}
  \dcase{\inj[!A_i]{i}{M}}{x_1}{N_1}{x_2}{N_2} \redone
  \Subst{N_i}{M}{x_i}
\end{gather*}
burada $M$,~$\delta$'nın; $N_i$ ise~$\delta_i$'nin ispat terimidir. Bu,
toplam tipleri için beta indirgeme kuralıdır. Burada
$\Subst{N_i}{M}{x_i}$, $N_i$ içindeki $x_i$ değişkeninin bütün serbest
geçişlerini~$M$ ile değiştirmek demektir.

Fonksiyon tipinin indirgemesi, \Intro\lif'ı izleyen \Elim\lif
dolambacının kaldırılmasına karşılık gelir:
\begin{gather*}
  \bottomAlignProof
  \AxiomC{$\Discharge{!A}{x}$}
  \RightLabel{$\delta$}
  \DeduceC{$!B$}
  \DischargeRule{$\Intro{\lif}$}{x}
  \UnaryInfC{$!A \lif !B$}
  \AxiomC{}
  \RightLabel{$\delta'$}
  \DeduceC{$!A$}
  \RightLabel{$\Elim{\lif}$}
  \BinaryInfC{$!B$}
  \DisplayProof\bottomAlignProof
  \quad
  \redone
  \quad
  \AxiomC{}
  \RightLabel{$\delta'$}
  \DeduceC{$!A$}
  \RightLabel{$\delta$}
  \DeduceC{$!B$}
  \DisplayProof
\end{gather*}
İspat terimleri için bu, sıradan beta indirgemesidir:
\begin{gather*}
  (\lambd[\typeof{x}{!A}][N]M)
  \redone \Subst{N}{M}{x}
\end{gather*}

!!{proof}s üzerindeki indirgeme dönüşümlerine ek olarak yer değiştirme
dönüşümlerini ele alabiliriz. Bunlar, \Elim{\lor} ile bitip ardından başka
bir !!{elimination} kuralı gelen (alt-)!!{proof}{}ına uygulanır; örneğin:
\begin{multline*}
  \bottomAlignProof
  \AxiomC{}
  \RightLabel{$\delta$}
  \DeduceC{$!A_i$}
  \RightLabel{$\Intro{\lor}$}
  \UnaryInfC{$!A_1 \lor !A_2$}
  \AxiomC{$\Discharge{!A_1}{x_1}$}
  \RightLabel{$\delta_1$}
  \DeduceC{$!C_1 \land !C_2$}
  \AxiomC{$\Discharge{!A_2}{x_2}$}
  \RightLabel{$\delta_2$}
  \DeduceC{$!C_1 \land !C_2$}
  \DischargeRule{$\Elim{\lor}$}{x_1,x_2}
  \TrinaryInfC{$!C_1 \land !C_2$}
  \RightLabel{\Elim{\land}[i]}
  \UnaryInfC{$!C_i$}
  \DisplayProof\bottomAlignProof
  \quad
  \redone
  \\
  \AxiomC{}
  \RightLabel{$\delta$}
  \DeduceC{$!A_i$}
  \RightLabel{$\Intro{\lor}$}
  \UnaryInfC{$!A_1 \lor !A_2$}
  \AxiomC{$\Discharge{!A_1}{x_1}$}
  \RightLabel{$\delta_1$}
  \DeduceC{$!C_1 \land !C_2$}
  \RightLabel{\Elim{\land}[i]}
  \UnaryInfC{$!C_i$}
  \AxiomC{$\Discharge{!A_2}{x_2}$}
  \RightLabel{$\delta_2$}
  \DeduceC{$!C_1 \land !C_2$}
  \RightLabel{\Elim{\land}[i]}
  \UnaryInfC{$!C_i$}
  \DischargeRule{$\Elim{\lor}$}{x_1,x_2}
  \TrinaryInfC{$!C_i$}
  \DisplayProof
\end{multline*}
$\delta$, $\delta_1$ ve $\delta_2$'nin ispat terimleri sırasıyla $M$,
$N_1$ ve $N_2$ ise ispat terimleri, yani lambda terimleri için karşılık
gelen indirgeme kuralı şöyledir:
\[
  \proj{i}{\dcase{M}{x_1}{N_1}{x_2}{N_2}} \redone
  \dcase{M}{x_1}{\proj{i}{N_1}}{x_2}{\proj{i}{N_2}}
\]

Elbette dolambaçları yalnız !!{proof}{}ın sonunda değil, dolambaçla biten
bir alt-!!{proof}{}a indirgeme dönüşümü uygulayarak !!{proof}{}ın içinde de
giderebiliriz. Benzer biçimde $\beta$-indirgemesi lambda terimlerinin alt
terimlerine uygulanır. $N$, $M$'nin bir alt terimi ve $N \redone N'$ ise
$M$ içindeki $N$ alt terimini~$N'$ ile değiştirip $M'$yi elde edebiliriz.
Bu durumda $M \redone M'$ de yazarız. $M = M_1 \redone M_2 \redone \dots
\redone M_k = M'$ biçiminde bir dizi varsa ($k=1$, yani $M=M'$ durumu
dâhil) $M \red M'$ yazarız.

Hiçbir alt-!!{proof}{}ına dönüşüm uygulanamayan bir !!a{proof}{}a
\emph{normal} denir. Benzer biçimde hiçbir (alt terimine) dönüşüm
uygulanamayan bir lambda terimine \emph{normal} denir. Normal bir lambda
terimine \emph{normal biçim} de denir.

\end{document}
